% ============================================================
%  preamble.tex — 論文「曲率制御付き Wasserstein 潜在拡散」
%  ML 会議（NeurIPS 型・単段組）に準じた体裁
% ============================================================

\usepackage{amsmath,amssymb,amsthm,mathtools}
\usepackage{tikz}
\usepackage{float}
\usepackage[dvipdfmx,hidelinks]{hyperref}
\usepackage{pxjahyper}
\usepackage{geometry}
\geometry{a4paper,margin=25mm}

% --- アルゴリズム環境（algorithm.sty に依存しない自前定義）---
\newfloat{algorithm}{tbp}{loa}
\floatname{algorithm}{アルゴリズム}
\newcounter{algstep}
\newenvironment{algsteps}{%
  \begin{list}{\arabic{algstep}:}{%
    \usecounter{algstep}%
    \setlength{\itemsep}{2pt}\setlength{\parsep}{0pt}\setlength{\topsep}{4pt}%
    \setlength{\leftmargin}{2.6em}\setlength{\labelwidth}{1.8em}%
    \setlength{\labelsep}{0.8em}%
  }}{\end{list}}

% --- 定理環境（論文体裁：装飾なし・節ごとの通し番号）---
\theoremstyle{plain}
\newtheorem{theorem}{定理}[section]
\newtheorem{proposition}[theorem]{命題}
\newtheorem{lemma}[theorem]{補題}
\newtheorem{corollary}[theorem]{系}

\theoremstyle{definition}
\newtheorem{definition}[theorem]{定義}
\newtheorem{fact}[theorem]{事実}

\theoremstyle{remark}
\newtheorem{remark}[theorem]{注意}

\renewcommand{\proofname}{証明}

% --- 記法 ---
\newcommand{\R}{\mathbb{R}}
\newcommand{\N}{\mathbb{N}}
\newcommand{\E}{\mathbb{E}}
\newcommand{\abs}[1]{\lvert #1 \rvert}
\newcommand{\norm}[1]{\lVert #1 \rVert}
\newcommand{\Pp}[1]{\mathcal{P}_{#1}}
\newcommand{\vol}{\operatorname{vol}}
\renewcommand{\d}{\mathrm{d}}

% --- 体裁 ---
\setlength{\parindent}{1zw}
\renewcommand{\abstractname}{要旨}
